% Generated draft: requires mathematical and source review.
\documentclass[11pt]{article}
\usepackage[margin=25mm]{geometry}
\usepackage{fontspec}
\setmainfont{Latin Modern Roman}
\usepackage{amsmath,amssymb,mathtools,unicode-math}
\setmathfont{Latin Modern Math}
\usepackage{enumitem,needspace,xurl,hyperref}
\setlength{\parindent}{0pt}
\setlength{\parskip}{5pt}
\setlength{\emergencystretch}{3em}
\newcommand{\N}{\mathbb{N}}
\newcommand{\Z}{\mathbb{Z}}
\newcommand{\Q}{\mathbb{Q}}
\newcommand{\R}{\mathbb{R}}
\newcommand{\C}{\mathbb{C}}
\newcommand{\F}{\mathbb{F}}
\newcommand{\A}{\mathbb{A}}
\newcommand{\PP}{\mathbb{P}}
\newcommand{\E}{\mathbb{E}}
\newcommand{\eps}{\varepsilon}
\newcommand{\dd}{\,\mathrm{d}}
\begin{document}
\section*{UnsolvedMath ID: 701532}
\textbf{Source problem code:} catalogue entry
\subsection{Definitions and mathematical statement}
Let \(p\) be a prime. A pro-\(p\) group is a topological group isomorphic, as a topological group, to an inverse limit of finite \(p\)-groups. For an abstract group \(H\), its commutator subgroup is
\[
[H,H]=\langle x^{-1}y^{-1}xy : x,y\in H\rangle,
\]
and \(H\) is perfect if \([H,H]=H\).

Construct a pro-\(p\) group \(G\) and a proper normal subgroup \(K\lhd G\), normal in the underlying abstract group, such that
\[
K\ne G,\qquad [G/K,G/K]=G/K.
\]
The quotient and its commutator subgroup are understood abstractly; no closedness assumption is imposed on \(K\), and no topological closure is taken in forming the commutator subgroup. Properness requires the perfect quotient to be nontrivial. The source does not specify whether \(p\) may be chosen as part of the example or whether an example is requested for every prescribed prime; this prime quantifier remains unresolved.
\subsection{Short English statement}
Construct a pro-\(p\) group with a nontrivial perfect abstract quotient by a proper normal subgroup, which is not required to be closed. The source leaves the intended quantifier over primes unspecified.
\subsection{Sources}
\begin{itemize}
\item Local UnsolvedMath registry, permanent ID \(701532\), source problem code \texttt{catalogue entry}, titled ``Kourovka Notebook Problem 18.104.'' The supplied registry labels the problem open; this label is not an independently verified mathematical or literature assessment.
\item \emph{The Kourovka Notebook (Unsolved Problems in Group Theory)}, original problem 18.104, proposed by John S. Wilson, 18th issue (2014); the supplied background also identifies the 21st issue edition (2026), arXiv:1401.0300v48. The supplied collection description names E. I. Khukhro and V. D. Mazurov as editors. These bibliographic details are supplied attribution and have not been independently verified.
\end{itemize}
\end{document}
